\documentclass[11pt,letterpaper]{article}

% ======= PACKAGES =======
\usepackage[utf8]{inputenc}
\usepackage{geometry}
\geometry{margin=1in}
\usepackage{setspace}
\setstretch{1.5}
\usepackage{fancyhdr}
\usepackage{xcolor}
\usepackage{titlesec}
\usepackage{lastpage}
\usepackage{hyperref}
\usepackage{mathpazo}
\usepackage{graphicx}
\usepackage{subcaption}
\linespread{1.5}

% ======= HEADER & FOOTER =======
\pagestyle{fancy}
\fancyhf{}

% Header (faded + smaller)
\fancyhead[L]{\footnotesize\textcolor{gray!90}{\textit{Data Management Plan}}}
\fancyhead[R]{\footnotesize\textcolor{gray!90}{\textit{Mostafa Rezaali}}}

% Footer (faded)
\fancyfoot[L]{\footnotesize\textcolor{gray!90}{\textit{NSF AGS-PRF -- University of Florida}}}
\fancyfoot[R]{\footnotesize\textcolor{gray!90}{Page \thepage}}

\renewcommand{\headrulewidth}{0.4pt}
\renewcommand{\footrulewidth}{0.4pt}

% ======= PARAGRAPH STYLE =======
\setlength{\parindent}{15pt}
\setlength{\parskip}{6pt}
\titlespacing*{\section}{0pt}{0.45em}{0.15em}
\titlespacing*{\subsection}{0pt}{0.35em}{0.1em}

% ======= DOCUMENT BEGINS =======
\begin{document}

% ======= TITLE BLOCK =======
\begin{center}
    \vspace*{-1.5cm}
    {\Large\bfseries Postdoctoral Fellowship Application} \\[0.2cm]
    {\small\textit{NSF Atmospheric and Geospace Sciences Postdoctoral Research Fellowship}} \\[0.05cm]
    {\small\textit{Department of Geography, University of Florida}} \\[0.1cm]
    \rule{0.6\textwidth}{0.4pt}
\end{center}
\vspace{-0.5cm}

\noindent
Mostafa Rezaali \\
Ph.D. Candidate, Climate Science (Geography) \\
University of Florida \\
mostafarezaali@gmail.com \quad $\mid$ \quad (352) 709-5350 \\[0.3cm]
{\large\bfseries Data Management Plan}\\[-0.15cm]
\rule{\textwidth}{0.4pt}

\section*{Products of Research}
The project will produce source code, trained-model configuration files, evaluation scripts, derived hindcast diagnostics, summary tables, selected nonrestricted figures, metadata, teaching materials, and manuscripts. The project will use third-party climate datasets such as PRISM and ERA5. Those source datasets will not be redistributed if license or provider terms require users to obtain them from the original sources.

\section*{Data and Code Formats}
Code will be released in plain-text Python and shell scripts with environment files where possible. Derived data products will be released in standard scientific formats such as NetCDF, NPZ, CSV, or GeoTIFF as appropriate, with README files describing variables, units, grids, time coverage, cross-validation splits, and known limitations. Figures and tables used in manuscripts will be preserved with scripts or notebooks sufficient to reproduce them from the derived products.

\section*{Access and Sharing}
Nonrestricted derived products and code will be shared through a public repository and archived with a persistent DOI through a repository such as Zenodo or an institutional archive. Large intermediate files and third-party raw data will be documented with scripts and instructions so that authorized users can reproduce them from original providers. Repository documentation will distinguish between raw third-party inputs, derived hindcast products, trained-model outputs, and manuscript-ready diagnostics.

\section*{Preservation}
Final code, metadata, manuscripts, teaching materials, and derived data needed to reproduce major figures and tables will be preserved for at least five years after the end of the award. Versioned releases will be used for manuscript-associated artifacts so that published results can be traced to a specific code and data snapshot.

\section*{Privacy, Security, and Ethical Considerations}
The project uses environmental and climate data, not human-subjects data. No personally identifiable information will be collected. Model limitations, uncertainty, and appropriate-use caveats will be documented to reduce the risk of overinterpreting experimental forecast products. Experimental forecasts will be presented as research outputs and not as official operational guidance.

\section*{Roles and Responsibilities}
The Fellow will be responsible for maintaining the repository, preparing metadata, documenting workflows, and archiving final products. Mentors will advise on scientific quality control, reproducible workflow design, and appropriate dissemination.
\end{document}
